% TOP_PROBLEM: {"id": 1176, "title": "Suslin's Problem", "category_id": 10, "category": {"id": 10, "name": "set_theory", "display_name": "Set Theory", "description": "Foundations of mathematics, infinite sets, and cardinality.", "slug": "set-theory", "order_index": 10, "created_at": "2026-07-31T15:26:25.671Z"}, "status": "open"}
% FIRST_POSTED: null
% ENTRY_KIND: statement_only
% Imported from: batch01.tex; UnsolvedMath ID 1176
\documentclass[11pt]{book}
\usepackage{fontspec}
\setmainfont{Latin Modern Roman}
\setsansfont{Latin Modern Sans}
\usepackage{amsmath,amssymb,mathtools}
\usepackage{unicode-math}
\setmathfont{Latin Modern Math}
\usepackage{xurl}
\usepackage[hidelinks]{hyperref}
\newcommand{\sref}[2]{\hyperlink{source:#1:#2}{[S#2]}}
\newcommand{\cataloguescope}[1]{\paragraph{Mathematical question.}#1}
\begin{document}
% UnsolvedMath ID 1176; source catalogue code SET-002
\setcounter{section}{1175}
\section{Suslin's Problem}\label{1176}
{\small\sffamily UnsolvedMath ID 1176 · Set Theory}
\subsection{Definitions and mathematical statement}

\paragraph{Shared notation.}$\mathbb N=\{1,2,\ldots\}$, and $\mathbb Z,\mathbb Q,\mathbb R,\mathbb C$ denote the integers, rationals, reals and complex numbers. Any other conventions are specified locally.


A nonempty linear order $(L,<)$ is dense if between any two distinct points there is a third point, has no endpoints if every point has a smaller and a larger point, and is Dedekind complete if every nonempty subset bounded above has a least upper bound in $L$. Its countable chain condition means that every family of pairwise disjoint nonempty open intervals $(a,b)=\{x:a<x<b\}$ is countable. An order isomorphism is a bijection preserving the order in both directions.
\paragraph{Statement recorded in the supplied source.}
Must every dense Dedekind-complete linear order without endpoints satisfying the countable chain condition be order-isomorphic to $(\mathbb R,<)$? \sref{1176}{2}
\subsection{Short English statement}
Do completeness, order density, lack of endpoints and countability of disjoint intervals characterize the ordered real line? This is the historical Suslin question; its answer is independent of ZFC, relative to the usual consistency assumption.
\subsection{Sources}
\begin{description}
\item[\hypertarget{source:1176:1}{S1}] ULAM.AI and the UnsolvedMath Contributors, \emph{UnsolvedMath: A Curated Collection of Open Mathematics Problems}, record 1176 (SET-002). CC BY 4.0. \url{https://huggingface.co/datasets/ulamai/UnsolvedMath}.
\item[\hypertarget{source:1176:2}{S2}] Chris Lambie-Hanson and Assaf Rinot, \emph{A forcing axiom deciding the generalized Souslin hypothesis}, arXiv:1708.06932v2 (2018). Section1.1, page1 and footnote1. \url{https://arxiv.org/pdf/1708.06932}.
\end{description}
\label{end:1176}
\end{document}
